\documentclass[11pt,letterpaper]{article}

\usepackage[margin=1in]{geometry}
\usepackage{setspace}
\onehalfspacing

\usepackage{microtype}
\usepackage{booktabs}
\usepackage{longtable}
\usepackage{array}
\usepackage{tabularx}
\usepackage{xcolor}
\usepackage{amsmath}
\usepackage[hyphens]{url}
\usepackage[colorlinks=true,linkcolor=blue,citecolor=blue,urlcolor=blue]{hyperref}
\usepackage{cleveref}
\usepackage[numbers,sort&compress]{natbib}

\title{
    \vspace{-1cm}
    \textbf{Update 12: USB4 RDMA Investigation} \\[0.5em]
    \large SoftRoCE over USB4 Networking for Distributed LLM Workloads
}

\author{
    Nathan Gopee \\
    \textit{MS Computer Science Candidate} \\
    SUNY New Paltz \\
    \texttt{gopeen1@newpaltz.edu}
}

\date{February 2026}

\begin{document}

\maketitle

\begin{abstract}
\noindent
This update investigates the feasibility of RDMA communication over USB4 networking between the Phase~1 cluster nodes (Cerberus and Chimera). USB4 offers a 20~Gbps point-to-point link that could supplement or replace the current 2.5~GbE RDMA path. We present hardware reconnaissance, driver analysis, and a gap assessment that identifies the requirements for USB4 RDMA deployment.
\end{abstract}

\section{Motivation}

Baseline SoftRoCE-based RDMA testing between Cerberus and Chimera over 2.5~GbE Realtek interfaces achieved 0.91~Gbps via \texttt{ib\_write\_bw} and 111.86~MB/s via UCX \texttt{tag\_bw}. While sufficient for protocol validation, the 2.5~GbE bottleneck limits practical tensor transfer throughput to $\sim$38\% of even the Realtek's theoretical line rate.

USB4 offers a potential 20~Gbps point-to-point link---an 8$\times$ bandwidth improvement over the current 2.5~GbE path. If SoftRoCE can be layered over a USB4 network interface, the resulting RDMA channel could approach the throughput of the 10~GbE NICs (Intel X710, Aquantia AQC107) without requiring direct NIC-to-NIC cabling.

\section{USB4 Architecture Overview}

USB4 is based on the Thunderbolt 3 protocol and provides:

\begin{itemize}
    \item \textbf{Tunneling architecture}: USB4 encapsulates multiple protocols (USB 3.x, DisplayPort, PCIe) within a single physical link.
    \item \textbf{IP networking}: The \texttt{thunderbolt-net} Linux kernel module creates a virtual Ethernet interface over the USB4/Thunderbolt tunnel, presenting a standard \texttt{thunderboltN} or \texttt{enpXsYuZuW} network device.
    \item \textbf{High MTU}: The \texttt{thunderbolt-net} driver supports MTU up to 65,520 bytes, dramatically reducing per-packet overhead compared to standard 1500-byte Ethernet.
    \item \textbf{Point-to-point}: USB4 networking is inherently point-to-point with no switching fabric, eliminating congestion from other traffic.
\end{itemize}

The RDMA stack layering for USB4 is:
\begin{center}
\begin{tabular}{@{}cl@{}}
\toprule
\textbf{Layer} & \textbf{Component} \\
\midrule
5 & Application (UCX \texttt{perftest}, NCCL) \\
4 & UCX / libibverbs \\
3 & SoftRoCE (\texttt{rdma\_rxe}) \\
2 & \texttt{thunderbolt-net} (virtual Ethernet) \\
1 & USB4 tunnel manager \\
0 & USB4 physical link (20 Gbps) \\
\bottomrule
\end{tabular}
\end{center}

\section{Hardware Reconnaissance}

\subsection{Cerberus: USB4 Controller Present}

\begin{table}[htbp]
\centering
\caption{Cerberus USB4/Thunderbolt hardware}
\label{tab:cerberus-usb4}
\begin{tabular}{@{}ll@{}}
\toprule
\textbf{Property} & \textbf{Value} \\
\midrule
Controller & ASMedia ASM4242 \\
PCI ID & \texttt{1b21:2426} (host), \texttt{1b21:2425} (downstream) \\
Kernel module & \texttt{thunderbolt} (loaded) \\
Sysfs domain & \texttt{domain0}, device \texttt{0-0} \\
USB4 device name & ASM4242 \\
Kernel & 6.17.0-14-generic \\
\bottomrule
\end{tabular}
\end{table}

Cerberus has a fully functional USB4 controller. The \texttt{thunderbolt} kernel module loaded automatically, and the ASM4242 controller registered as \texttt{domain0} in \texttt{/sys/bus/thunderbolt/devices/}. The controller supports USB4 tunneling including IP networking via \texttt{thunderbolt-net}.

\subsection{Chimera: USB 3.2 Only}

\begin{table}[htbp]
\centering
\caption{Chimera USB controller hardware}
\label{tab:chimera-usb4}
\begin{tabular}{@{}ll@{}}
\toprule
\textbf{Property} & \textbf{Value} \\
\midrule
Controller & ASMedia ASM3242 \\
PCI ID & \texttt{1b21:3242} \\
Protocol & USB 3.2 Gen 2$\times$2 (20 Gbps) \\
Kernel driver & \texttt{xhci\_hcd} \\
Thunderbolt & \textbf{Not detected} \\
Motherboard & ASRock TRX40 Creator (TR 3960X) \\
TB3 header & Present on board (unpopulated) \\
\bottomrule
\end{tabular}
\end{table}

Chimera's ASM3242 is a USB 3.2 host controller---not USB4/Thunderbolt. The \texttt{thunderbolt} kernel module loaded without error but detected no devices in \texttt{/sys/bus/thunderbolt/devices/}. The ASRock TRX40 Creator motherboard has an onboard Thunderbolt 3 header (Intel Titan Ridge compatible), but this requires an add-in card (AIC) that is not currently installed.

\subsection{Protocol Mismatch Analysis}

USB4 networking via \texttt{thunderbolt-net} requires a Thunderbolt/USB4 tunnel between both endpoints. The tunnel negotiation fails when:

\begin{enumerate}
    \item The remote device does not have a USB4/Thunderbolt controller (current situation).
    \item The cable does not support USB4 signaling (passive USB-C cables may lack Thunderbolt support).
    \item The Thunderbolt security policy rejects the remote device (managed by the \texttt{bolt} daemon).
\end{enumerate}

In our case, condition (1) is the blocking issue: Cerberus has USB4 (ASM4242) but Chimera has only USB 3.2 (ASM3242). No Thunderbolt tunnel can be established, and therefore no \texttt{thunderbolt-net} interface will appear.

\section{Attempted Configuration}

Despite the hardware mismatch, the following steps were attempted to verify the diagnosis:

\begin{enumerate}
    \item \textbf{Cerberus}: The \texttt{thunderbolt} and \texttt{thunderbolt-net} kernel modules were loaded (the former was already auto-loaded). Inspecting \texttt{/sys/bus/thunderbolt/devices/} confirmed the local ASM4242 controller (\texttt{domain0}, device \texttt{0-0}) but no remote peer device.

    \item \textbf{Chimera}: Both modules loaded without error, but \texttt{/sys/bus/thunderbolt/devices/} was empty---confirming that the ASM3242 does not register as a Thunderbolt device.

    \item \textbf{USB gadget alternative}: Inspecting \texttt{/sys/class/udc/} showed no USB Device Controllers available. The \texttt{g\_ncm} (USB NCM gadget) module loaded but could not bind without a UDC. Desktop xHCI controllers support host mode only; OTG or device-capable USB controllers are required. The \texttt{dummy\_hcd} virtual UDC module was not available on kernel 6.17.0-14.
\end{enumerate}

\section{Expected Performance (Projected)}

Based on USB4 20~Gbps specifications and SoftRoCE overhead characteristics:

\begin{table}[htbp]
\centering
\caption{Projected USB4 RDMA performance vs.\ current 2.5~GbE baseline}
\label{tab:usb4-projected}
\begin{tabular}{@{}lrrl@{}}
\toprule
\textbf{Metric} & \textbf{2.5 GbE (actual)} & \textbf{USB4 20G (projected)} & \textbf{Improvement} \\
\midrule
Raw link rate & 2.5 Gbps & 20 Gbps & 8$\times$ \\
\texttt{ib\_write\_bw} & 0.91 Gbps & 10--15 Gbps & 11--16$\times$ \\
UCX \texttt{tag\_bw} (1 MB) & 0.89 Gbps & 8--12 Gbps & 9--13$\times$ \\
Latency (1 byte) & $>$100 $\mu$s & 15--50 $\mu$s & 2--7$\times$ \\
Max MTU & 1500 B & 65,520 B & 44$\times$ \\
\bottomrule
\end{tabular}
\end{table}

\paragraph{Key assumptions.}
\begin{itemize}
    \item MTU 65,520 bytes reduces per-packet overhead by $\sim$44$\times$ compared to standard Ethernet, significantly benefiting large tensor transfers.
    \item SoftRoCE CPU overhead remains the primary bottleneck; both machines have high-core-count CPUs (48 logical cores on Cerberus, 48 on Chimera) to absorb this.
    \item The USB4 link is point-to-point with no competing traffic, unlike the shared 2.5~GbE LAN.
    \item GPUDirect RDMA is not possible over SoftRoCE regardless of link speed; all transfers must transit system memory.
\end{itemize}

\section{Resolution Path}

\subsection{Option A: Thunderbolt 3 AIC for Chimera (Recommended)}

The ASRock TRX40 Creator has a dedicated Thunderbolt 3 header. Installing a compatible Thunderbolt 3 Add-In Card (e.g., ASRock Thunderbolt 3 AIC R2.0 or GC-Titan Ridge) would:

\begin{enumerate}
    \item Add an Intel Titan Ridge (JHL7540) Thunderbolt 3 controller to Chimera.
    \item Enable USB4/Thunderbolt tunnel negotiation with Cerberus's ASM4242.
    \item Create a \texttt{thunderbolt-net} interface on both machines.
    \item Allow SoftRoCE to bind to the high-speed tunnel interface.
\end{enumerate}

\textbf{Estimated cost:} \$30--60 for a Thunderbolt 3 AIC.

\subsection{Option B: Direct 10 GbE Cabling}

Both machines have 10~GbE NICs (Intel X710 on Cerberus, Aquantia AQC107 on Chimera). A direct DAC/SFP+ cable between them would provide:

\begin{itemize}
    \item 10~Gbps full-duplex bandwidth.
    \item Native iWARP RDMA on the X710 side (Aquantia side still requires SoftRoCE).
    \item No additional hardware purchase required if a 10GbE DAC cable is available.
\end{itemize}

\subsection{Option C: USB 3.2 Gadget Networking}

If a USB OTG-capable adapter or bridge cable is available, USB 3.2 Gen 2$\times$2 networking could provide 20~Gbps. However:
\begin{itemize}
    \item Requires either a USB bridge cable (e.g., Plugable USB3-TRAN) or USB OTG support.
    \item Standard desktop USB controllers do not support gadget mode.
    \item Performance would be limited by the USB 3.2 protocol overhead.
\end{itemize}

\section{Conclusions}

\begin{enumerate}
    \item \textbf{USB4 RDMA is architecturally sound}: SoftRoCE over \texttt{thunderbolt-net} is a valid approach that requires no kernel patches---only standard modules (\texttt{thunderbolt}, \texttt{thunderbolt-net}, \texttt{rdma\_rxe}).

    \item \textbf{Hardware gap identified}: Chimera lacks a USB4/Thunderbolt controller. The ASM3242 USB 3.2 controller cannot establish Thunderbolt tunnels. The TRX40 Creator's onboard TB3 header provides a clear upgrade path via an inexpensive AIC.

    \item \textbf{Projected 8--16$\times$ improvement}: USB4 20~Gbps with MTU 65,520 should deliver 10--15~Gbps RDMA bandwidth, compared to 0.91~Gbps on the current 2.5~GbE path.

    \item \textbf{Next step}: Acquire a Thunderbolt 3 AIC for Chimera, install it on the TRX40 Creator's TB3 header, then repeat the SoftRoCE/UCX test suite over the \texttt{thunderbolt-net} interface.
\end{enumerate}

\section{Deployment Procedure (For Future Deployment)}
\label{sec:usb4-setup}

Once both endpoints have compatible USB4/Thunderbolt controllers, the deployment procedure is:

\begin{enumerate}
    \item \textbf{Module loading}: Load \texttt{thunderbolt}, \texttt{thunderbolt-net}, and \texttt{rdma\_rxe} kernel modules on both machines.

    \item \textbf{Thunderbolt authorization}: Install the \texttt{bolt} daemon and enroll all discovered Thunderbolt devices with \texttt{boltctl enroll} using an \texttt{auto} security policy. This authorizes the peer controller for tunnel establishment.

    \item \textbf{Interface discovery}: After a brief settling period (${\sim}3$~s), scan \texttt{/sys/class/net/} for the \texttt{thunderbolt-net}-backed interface by following each device's driver symlink.

    \item \textbf{Network configuration}: Set MTU to 65,520 bytes (the \texttt{thunderbolt-net} driver maximum), assign point-to-point IP addresses (e.g., 10.10.10.1/24 and 10.10.10.2/24), and bring the interface up.

    \item \textbf{SoftRoCE binding}: Create an \texttt{rxe0} RDMA device bound to the Thunderbolt network interface via \texttt{rdma link add rxe0 type rxe netdev <iface>}.

    \item \textbf{Validation}: Verify connectivity with \texttt{ping}, then run \texttt{ib\_write\_bw -d rxe0 --report\_gbits} and \texttt{ucx\_perftest} with \texttt{UCX\_NET\_DEVICES} pinned to the Thunderbolt interface.
\end{enumerate}

% No citations in this update

\end{document}
